\chapter{Analysis from the files}
\label{ch:analysis}

\chaptergoals{10 of 11 in Part~II.}
{\cref{ch:aegir,ch:digitisation}: the simulated and digitised files; the Python of \cref{ch:team}.}
{select HNL decays from the files, weight them correctly, turn efficiencies into yields, and measure momentum from the bending.}

There is no analysis framework to learn in this chapter. The files are
RNTuples of plain structs, so a Python script with PyROOT can read them
directly, and that is what we do.

\section{Reading the files}

The pattern from \cref{ch:generator} is all you need: open a reader on the
RNTuple, create a typed \emph{view} per field, and call the view with an entry
number to get that entry's value as a C++ object. The scripts in
\file{code/analysis/} share three small helpers in \file{shipfiles.py}: a
class that keeps readers and views together, the detector numbering, and the
positions of the decay volume and the straw stations.

\codefile[linerange={36-63}]{Python}{code/analysis/shipfiles.py}

\begin{pitfall}
Two habits of PyROOT with RNTuples are built into these helpers. A view
returns a reference into its own buffer, which the next call overwrites: copy
out what you need before reading the next entry. And tearing down readers
when Python exits can crash; the scripts therefore keep every reader alive and
leave through \texttt{os.\_exit} after flushing their output
(\texttt{sf.finish()}).
\end{pitfall}

\section{Which file for which question}

\begin{center}
\small
\begin{tabularx}{\linewidth}{X l}
\toprule
Question & File \\
\midrule
generator truth: HNL mass, decay position, daughter momenta & input or simulation \\
per event: which detectors did the daughters reach, was the veto silent & simulation \\
weighted efficiencies and yields & simulation (carries the weights) \\
detector response over all hits: resolutions, time distributions & digitisation \\
\bottomrule
\end{tabularx}
\end{center}

The reason for the split is the shape of \pkg{Shannon}'s output
(\cref{sec:digilayout}): a stream of time-stamped hits, as the detector will
deliver, with no simulated-event number and no weights. Questions about one
simulated event are therefore answered from the simulation file, using
simulated hits as a stand-in for digitised ones. For the coarse station-count exercise, the position smearing is small.
This does not establish equivalence for reconstructed tracks, vertex cuts,
timing or a realistic detector response.

\section{The HNL analysis}

\file{code/analysis/hnl\_analysis.py} answers, event by event:

\begin{enumerate}
  \item What does the generator record say? It requires a pion and a muon of opposite electric charge
        with status~1, a shared status-2 mother and a common vertex, and fills the $\pi\mu$ invariant mass, the
        decay position and the opening angle.
  \item Where did the daughters go? It finds them among the Geant4 particles
        (\field{parentId}~$=0$), and for each counts the straw stations in which
        it left hits.
  \item Was the veto silent? It counts UBT hits and sums the SBT energy
        deposit.
  \item How much did the magnet bend each daughter? From the mean $x$ of its
        hits in stations 1--2 and 3--4 it computes the change of slope
        $\Delta\theta_x$ across the magnet.
\end{enumerate}

It then applies these as a sequence of cuts, weighting every event with the
EventCalc weight from \texttt{event\_header}, and prints the cut flow.

\begin{exercise}{Run the analysis}
From your working directory:
\begin{shell}
python $CODE/analysis/hnl_analysis.py hnl_sim.root
\end{shell}
Look at the printed cut flow and at the PDFs it writes
(\file{hnl\_analysis\_*.pdf}).
\begin{enumerate}
  \item \file{\_mass.pdf}: why is the peak so narrow, and where is it?
  \item \file{\_decay\_z.pdf}: describe the shape. How would it change for a
        much shorter $c\tau$?
  \item Which cut removes most of the weighted events? Look at
        \file{\_n\_stations.pdf} and \file{\_opening.pdf} to find out why.
\end{enumerate}
\end{exercise}

\begin{checkpoint}
A cut flow of seven lines whose first line counts all your simulated events,
and eight PDFs.
\end{checkpoint}

\begin{solution}
\begin{enumerate}
  \item The mass is computed from the generator four-vectors, which conserve
        energy and momentum exactly: the peak is a spike at the HNL mass of
        the sample. The width a real analysis sees comes from the momentum
        resolution of reconstruction.
  \item For laboratory decay length $\beta\gamma c\tau$ large compared with the volume length, the
        decay density is approximately uniform along it; the slope is the
        exponential decay law, steeper for shorter $c\tau$, which piles decays
        up at the upstream end. The EventCalc weights carry the
        normalisation, which is why the histogram is filled with them.
  \item Usually the station requirement: daughters at large angle leave the
        spectrometer acceptance, and soft ones are swept out by the magnet
        before stations 3 and 4.
\end{enumerate}
\end{solution}

\begin{pitfall}
The cut labelled ``decay in nominal z range'' checks only the longitudinal
coordinate, not the helium volume or a fiducial margin. The UBT and SBT
criteria also inherit the broad sensitive-volume selection. A full analysis
needs a three-dimensional geometry check, active-material hit selection,
reconstruction and background samples before quoting a sensitivity.
\end{pitfall}

\section{The veto}

The SBT surrounds the decay volume to reject background from particles
entering it from outside. Signal daughters should not touch it, but they can:
a daughter emitted at a large angle, or a secondary from an interaction in the
vessel or the first station, deposits energy in the LAB cells.

\begin{exercise}{What does a tight veto cost?}
\file{\_sbt\_edep.pdf} shows the SBT energy deposit per event. Rerun with
different thresholds and record the signal efficiency of the
``SBT below threshold'' line:
\begin{shell}
for t in 0 0.010 0.045 0.100; do
  python $CODE/analysis/hnl_analysis.py hnl_sim.root --sbt-threshold $t | grep SBT
done
\end{shell}
A threshold costs signal efficiency here, and buys background rejection in a
background sample. Why is the background side of this trade-off not
something you can study with this sample?
\end{exercise}

\begin{solution}
The sample contains only signal. The rejection of a given threshold has to be
measured on muon-induced and neutrino-induced background samples, simulated
with the \texttt{fixed\_target\_source} or \texttt{genie\_reader\_source}
plugins of \pkg{aegir}, which is a natural follow-up project.
\end{solution}

\section{Measuring momentum with the magnet}

A charged particle crossing a dipole field changes direction by
\[
  \Delta\theta \approx \frac{0.3\,\int B\,\mathrm{d}l\;[\mathrm{T\,m}]}{p\;[\mathrm{GeV}]},
\]
for a singly charged, forward track in the small-angle limit. Thus
$p\,|\Delta\theta_x|$ is approximately constant only for similar sampled
field integrals and negligible energy loss and scattering.
\file{\_p\_kick.pdf} plots it with the true momentum $p$ and the
$\Delta\theta_x$ measured from the simulated straw hits.

\begin{exercise}{The field integral from the hits}
Read the position of the peak in \file{\_p\_kick.pdf} and convert it to
$\int B\,\mathrm{d}l$. Compare with your estimate from the field map in
\cref{sec:fieldservice}. If you simulated without field in \cref{ch:aegir},
run the analysis on that file too: where is the peak now?
\end{exercise}

\begin{solution}
The peak is at $0.3\int B_y\,\mathrm{d}z$ in GeV\,rad, e.g.\ \tbd{value for
the tutorial map}; the distribution is broadened by varying paths, scattering and energy
loss; primary-track matching excludes secondary hits from this estimator. Without field
the peak sits at zero, broadened by scattering and the finite sampling of the trajectory. Turning this
around, $p = 0.3\int B\,\mathrm{d}l / |\Delta\theta|$ is the momentum
measurement reconstruction makes, with a fitted track instead of mean hit
positions.
\end{solution}

\section{Yields}

The weights carry the probability of each decay; the absolute normalisation is
EventCalc's expected number of decays $N_\mathrm{exp}$
(\cref{ch:generator}). Provided its exposure, channel selection and sampled phase space match
the input (with no omitted acceptance), it converts each weighted cut-flow
fraction into a teaching yield estimate:
\begin{shell}
python $CODE/analysis/hnl_analysis.py hnl_sim.root --n-expected <N_exp>
\end{shell}

\begin{exercise}{How many would SHiP see?}
Take $N_\mathrm{exp}$ for the sample from \tbd{where the converter or the
organisers provide it} and compute the expected number of selected
$\HNL\to\pi\mu$ decays. Which of the selection steps here would a real analysis
replace, and with what?
\end{exercise}

\begin{solution}
Truth-level and simulated-hit requirements stand in for reconstruction: the
station requirement becomes ``two reconstructed tracks with a good fit'', the
decay-volume cut a fiducial cut on the \emph{reconstructed} vertex, and the
analysis adds a cut on the impact parameter of the HNL candidate with respect
to the target, the invariant mass, muon identification and timing coincidence.
All of these need \pkg{Trout}.
\end{solution}

\section{Going further}

\begin{exercise}[colback=white]{Optional: a vertex without reconstruction}
Upstream of the magnet, straight lines approximate the daughter trajectories;
material scattering and energy loss still matter. For each daughter,
take the mean position of its simulated straw hits in stations 1 and 2,
extrapolate the line back into the decay volume, and find the point of
closest approach of the two lines. Histogram its distance from the true decay
vertex in $x$ and $z$. Why is the resolution so much worse in $z$ than in
$x$, and how does it depend on the opening angle?
\end{exercise}

\begin{solution}
A solution is in \file{code/analysis/solutions/vertex\_from\_hits.py}. The
lines meet at a small angle, so an error in their transverse position becomes
an error in $z$ larger by roughly the inverse opening angle; small-angle
decays (light, fast HNLs) have the worst $z$ resolution.
\end{solution}

\begin{exercise}[colback=white]{Optional: from a hit stream back to events}
Using the digitised file, histogram the times of all straw hits. Could you
separate the simulated events by time? What would the digitisation have to
record so that a Monte Carlo study can connect digitised hits to the
simulated event they came from? Bring your answer to the closing discussion.
\end{exercise}
